% Revised from the author-supplied schubmult.tex snapshot, 1 November 2026.
% The original manuscript is not modified.
\documentclass{amsart}
\usepackage{amsmath,amssymb,amsthm,mathtools}
\usepackage[margin=1in]{geometry}
\usepackage[hidelinks]{hyperref}
\hypersetup{
 pdftitle={Input-parameter bounds for (quantum) single and double Schubert multiplication by elementary expansions},
 pdfauthor={Matthew J. Samuel}
}
\newtheorem{theorem}{Theorem}[section]
\newtheorem{proposition}[theorem]{Proposition}
\newtheorem{lemma}[theorem]{Lemma}
\newtheorem{corollary}[theorem]{Corollary}
\theoremstyle{definition}
\newtheorem{definition}[theorem]{Definition}
\newtheorem{method}[theorem]{Method}
\newtheorem{example}[theorem]{Example}
\theoremstyle{remark}
\newtheorem{remark}[theorem]{Remark}
\numberwithin{equation}{section}
\newcommand{\Sch}{\mathfrak S}
\newcommand{\code}{\mathfrak c}
\newcommand{\id}{\mathrm{id}}
\newcommand{\Des}{\mathrm{Des}}
\newcommand{\fall}[2]{(#1)_{\underline{#2}}}
\newcommand{\alg}{\texttt{schubmult}}

\title[Input-parameter bounds for Schubert multiplication]{Input-parameter bounds for (quantum) single and double Schubert multiplication by elementary expansions}
\author{Matthew J. Samuel}
\email{matthematics@gmail.com}
\date{November 1, 2026}
\subjclass[2020]{05E05, 05E14, 14N15, 68Q25}
\keywords{Schubert polynomial, double Schubert polynomial, Pieri rule,
Lehmer code, quantum cohomology, arithmetic complexity}

\begin{document}
\begin{abstract}
We bound the complexity of the elementary-expansion algorithm of \alg,
which computes entire products of Schubert polynomials, in parameters
read from the input: the ambient permutation size $n$, the lengths
$r_u,r_v$ of the trimmed Lehmer codes, and the total inversion number $d$.
A single analysis covers ordinary Schubert multiplication, double
Schubert multiplication, the \emph{triple} product
$\Sch_u(x;y)\Sch_v(x;z)$ with independent coefficient alphabets, and
the quantum version of each.
With $r=\max(r_u,r_v)$ and $h=r_v$, the full triple product uses at most
$O(2^{2r}n^{2r+2h+4})$ elementary and coefficient-ring operations.
This is $2^{O(n\log n)}$ in general and polynomial in $n$ for fixed code
lengths, and it is the same bound we obtain for ordinary products.
When $u$ and $v$ are Grassmannian with descent at $k$, every $u$-side
state of the classical computation is $k$-Grassmannian.

This uniformity is a feature of the method.
It does not follow from specialization, and reducing triple
coefficients to ordinary ones monomial by monomial would multiply the
cost by the size of the expanded output.
The factorial elementary expansion of $\Sch_v(x;z)$ is indexed by
exactly the same column paths as the ordinary elementary expansion of
$\Sch_v(x)$.
The double Pieri rule is indexed by exactly the same chains as the
ordinary one, each carrying a single factorial elementary coefficient.
The two extra alphabets therefore enlarge neither the state space nor
the label graphs; they cost $O(h^2)$ ring operations per matched pair.

For the quantum triple product, and hence for equivariant and ordinary
quantum Littlewood--Richardson coefficients of complete flag varieties,
the bound has the same form. It is again $2^{O(n\log n)}$ in general and
polynomial in $n$ when $r_u$, $r_v$, and the number $L$ of nonzero
entries of the code of $v^{-1}$ are fixed.
The key observations are that the Fomin--Gelfand--Postnikov quantization
commutes with the filtered expansion when column heights are distinct,
and that each quantum Pieri step can raise the last descent by at most
one, at a cost in $q$-degree.

Coefficients are represented by arithmetic circuits; no claim about
expanded monomial output, symbolic zero testing, or polynomiality in
the binary encoding of code entries is made.
We compare these full-product bounds with Pak and Slonim's
bounded-range algorithm for a single ordinary Schubert coefficient.
\end{abstract}
\maketitle

\section{Introduction}

For permutations $u,v$ with finite support, consider the expansion
\begin{equation}\label{eq:product}
 \Sch_u(x;y)\Sch_v(x;z)
   =\sum_w c_{u,v}^w(y;z)\Sch_w(x;y).
\end{equation}
Here $x,y,z$ are independent infinite alphabets, and the coefficients belong
to $\mathbb Z[y,z]$.
We call \eqref{eq:product} the \emph{triple} product.
At $z=y$ its coefficients are the equivariant Schubert structure
constants of complete flag varieties.
At $y=z=0$ they are the ordinary Schubert structure constants, which
are nonnegative integers.
For $y\ne z$ the coefficients are more general than both.
The problem in this paper is to compute all the coefficients in
\eqref{eq:product}, not just a coefficient specified by a third permutation.

The \alg\ package \cite{schubmult} computes such products by matching two
computations. An elementary-symmetric expansion of the second factor
provides a graph of signed labels. Pieri's rule multiplies the first factor
by each elementary factor. Dynamic programming combines these operations
without expanding the final coefficients into monomials.
The expansion formulas and correctness of this architecture are established
in \cite[Section~4 and Section~7]{samuelquantum}.
Here we give an input-only complexity analysis, with the non-quantum
algorithm specified precisely enough to account for every search node.
We use ``classical'' only to mean ``non-quantum''; it includes the
double and triple products.

The analysis does not use the number of nonzero output coefficients,
a worst-case output size, or the support of a different product as a
complexity parameter.
We use the distinct-position column paths of
\cite[Definition~4.4 and Theorem~4.5]{samuelquantum};
their restriction is justified by the cancellation theorem proved there,
not inferred from an endpoint count.
The arithmetic-circuit model and the choice of elementary-weight evaluator
are distinguished from literal symbolic implementation costs in
Section~\ref{sec:implementation}.

\subsection{Parameters and main bound}

Permutations are composed as functions, with $(ab)(i)=a(b(i))$.
Write
\[
 \code_i(w)=\#\{j>i:w(j)<w(i)\},\qquad
 \ell(w)=\sum_i\code_i(w),\qquad
 r_w=\max\{i:\code_i(w)>0\},
\]
where $r_{\id}=0$. Thus $r_w$ is the length of the Lehmer code
\emph{after deleting trailing zeros}, not the number of nonzero entries.
It is also the last descent position of $w$.
We take $n\geq2$ to be the least integer for which $u,v\in S_n$,
apart from this harmless lower bound of $2$.
It can be recovered from the input codes:
\[
 n=\max\bigl(\{2\}\cup\{i+\code_i(u):\code_i(u)>0\}
                     \cup\{i+\code_i(v):\code_i(v)>0\}\bigr).
\]
Put $r=\max(r_u,r_v)$, $h=r_v$, and $d=\ell(u)+\ell(v)$.
The case $v=\id$ requires no multiplication and will be set aside.

Section~\ref{sec:completion} constructs a dominant completion with column
heights $\theta_1\geq\cdots\geq\theta_L>0$.
It satisfies $\theta_1=h$ and $L\leq n-1$.
Define the following quantities, all determined or bounded by the input:
\begin{equation}\label{eq:parameters}
 \begin{split}
 M&=n+L,\\
 U&=\min\left\{\fall{M}{r},\binom{d+r}{r}\right\},\\
 V&=\fall{n}{h},\qquad
 Q_h(n)=\sum_{b=0}^{h}\fall{n-1}{b}.
 \end{split}
\end{equation}
Here $\fall{a}{b}=a!/(a-b)!$ and $\fall{a}{0}=1$.
The quantities $U$ and $V$ are explicit bounds on state spaces,
not observed support sizes.

\begin{theorem}[Input-parameter bound]\label{thm:main}
For $v\ne\id$, Method~\ref{method:product} computes the entire triple
product \eqref{eq:product}, with coefficient arithmetic circuits, in
\begin{equation}\label{eq:main}
 O\left(n^3+L M^3 U^2 V Q_h(n)\right)
\end{equation}
elementary and coefficient-ring operations.
The same bound, with the same proof, holds for the double ($z=y$) and
ordinary ($y=z=0$) products; see Proposition~\ref{prop:uniform}.
In particular, the bound is
\begin{equation}\label{eq:coarse}
 O\left(2^{2r}n^{2r+2h+4}\right)
   =2^{O(n\log n)}.
\end{equation}
For fixed $r_u,r_v$ it is polynomial in $n$.
If $v$ is dominant, the bound improves to
$O(n^3+L M^3 U^2)$.
\end{theorem}

The proof is in Section~\ref{sec:complexity}.
The leading constant in the general exponent is not claimed to be
optimal: \eqref{eq:coarse} gives at most $4n\log_2 n+O(n)$.
In particular, we do not retain a more precise exponent without
establishing every support and branching estimate needed for it.

\subsection{Quantum multiplication}

The same architecture computes the triple quantum product
\eqref{eq:qproduct} of quantum double Schubert polynomials, whose
specializations are the equivariant and ordinary quantum
Littlewood--Richardson coefficients of complete flag varieties.
Two further input parameters enter.
Let $L$ be the number of nonzero entries of $\code(v^{-1})$; it is
the number of columns of the completion above.
Let $\tilde h$ be the largest part of the strict completion
(Method~\ref{method:strict}).
It is computed from $v$ in $O(n^3)$ operations and, by
Lemma~\ref{lem:strict}, satisfies
\[
 \max(r_v,L)\leq\tilde h\leq\min(n-1,\ r_v+L-1).
\]
Put $\rho=\max(r_u,\tilde h)$, $\widetilde V=\fall{n}{\tilde h}$, and
define $R_q\leq2\rho$ and $U_q$ as in Section~\ref{sec:quantum}.

\begin{theorem}[Quantum input-parameter bound]\label{thm:quantum}
For $v\neq\id$, Method~\ref{method:qproduct} computes every
coefficient $c_{u,v}^{w,\delta}(y;z)$ in \eqref{eq:qproduct}, as
arithmetic circuits over $\mathbb Z[q][y,z]$, in
\begin{equation}\label{eq:qmain}
 O\left(n^3+LM^3U_q^2\widetilde VQ_{\tilde h}(n)\right)
\end{equation}
elementary and coefficient-ring operations.
Every $w$ that occurs lies in $S_{n+L}$, and every $q^\delta$
that occurs involves only $q_1,\ldots,q_{n+L-1}$; hence the
coefficients at $z=y$ are the equivariant quantum
Littlewood--Richardson coefficients of $Fl_N$ for every $N\geq n+L$.
The bound is
\begin{equation}\label{eq:qcoarse}
 O\left(2^{2R_q}n^{2R_q+2\tilde h+4}\right)
 \subseteq O\left(2^{4\rho}n^{6\rho+4}\right)
 =2^{O(n\log n)},
\end{equation}
and it is polynomial in $n$ for fixed $r_u,r_v,L$.
The same bounds hold for the equivariant ($z=y$) and ordinary
($y=z=0$) quantum products, again with the same state space.
\end{theorem}

Unlike Theorem~\ref{thm:main}, the polynomial case requires $L$ to be
fixed. For $\Sch_v=x_1^m$ one has $r_v=1$ but $\tilde h=m$;
see Section~\ref{sec:quantum}.

\subsection{One bound for ordinary, double, triple, and quantum products}
\label{sec:uniform}

Theorems~\ref{thm:main} and \ref{thm:quantum} are stated for the most
general products, but the bounds are not weakened by that generality.
The ordinary product is not analyzed separately and then lifted.
Instead, the triple computation and its ordinary specialization build
the same label graphs, make Pieri searches over the same chains, and
are bounded by the same state counts.
The two coefficient alphabets add only the cost of evaluating one
factorial elementary polynomial per matched pair
(Proposition~\ref{prop:uniform}).

This depends on two properties of the method.
\begin{enumerate}
\item The factorial elementary expansion of $\Sch_v(x;z)$
(Theorem~\ref{thm:elementary}) is indexed by the same filtered column
paths as the ordinary elementary expansion of $\Sch_v(x)$.
The alphabet $z$ enters only through the weight attached to each path.
\item The double Pieri rule \eqref{eq:pieri} is indexed by the same
canonical chains as Sottile's ordinary Pieri rule.
Each chain carries a single factorial elementary polynomial in a fixed
alphabet, rather than a sum over additional combinatorial data.
\end{enumerate}
The quantum Pieri rule \eqref{eq:qpieri} has the same feature, with
a $q$-monomial attached to each chain.
Its chains follow the quantum Pieri formula of Postnikov
\cite{postnikov1999quantum} and Ciocan-Fontanine
\cite{ciocan1999partial} rather than Sottile's rule.
Quantum covers add paths: they need not increase length, and the
strict completion can be taller than the ordinary one.
The extra paths are absorbed by the same counting, with $\tilde h$ and
$R_q$ in place of $h$ and $r$, so the asymptotic bound is essentially
unchanged (Theorem~\ref{thm:quantum}).

Neither property is a formal consequence of specialization.
Ordinary structure constants do determine the triple coefficients, but
the resulting reduction does not give a uniform bound.
By Hwang's coefficient lifting \cite[Theorem~1.2]{hwang2026lifting},
the coefficient of $y^\alpha$ in $c_{u,v}^w(y;z)$ is the single
coefficient $c_{\hat u,\hat v}^{\hat w}(0;z)$ for an explicit
block-lifted triple $(\hat u,\hat v,\hat w)$ in $S_{n+|\alpha|}$, where
$|\alpha|\leq\ell(u)+\ell(v)-\ell(w)$.
(Hwang expands in the Schubert basis of the second factor's alphabet,
so the two factors of \eqref{eq:product} are exchanged in his
notation; in ours the lifted alphabet $y$ is that of the output basis.)
The Cauchy formula then writes the coefficient of each monomial
$y^\alpha z^\beta$ as $\pm$ a sum of products of ordinary structure
constants \cite[Section~5.2]{hwang2026lifting}.
Thus an algorithm for $c_{u,v}^w$ does yield $c_{u,v}^w(y;z)$, but not
parsimoniously.
The reduction runs separately for each monomial and each Cauchy term.
Each run uses permutations whose rank exceeds $n$ by up to the degree of
the monomial.
Its cost is therefore multiplied by the expanded size of the triple
coefficients, and its input parameters grow.
Here, by contrast, the double and triple coefficients come at no extra
asymptotic cost because they are produced as arithmetic circuits along
the same paths as their ordinary specializations.
The circuit representation is essential to this statement, and we make
no comparable claim for expanded monomial output.

\subsection{Prior work and the encoding distinction}\label{sec:prior}

Pieri formulas for ordinary Schubert polynomials were established by
Sottile \cite{sottile}; we proved the double Pieri formula used here in
\cite[Theorem~7.1]{samuelmolev}.
The Cauchy identity and dominant factorization are standard
\cite{notes}.
We studied the triple coefficients in this generality in
\cite{samuelmolev}.
Hwang \cite{hwang2026lifting} recently showed that each of their
monomial coefficients is determined by ordinary structure constants
through an explicit block lift; Section~\ref{sec:uniform} compares that
reduction with the present method.
In earlier work \cite{samuelquantum}, we established the ordinary-double
expansion \cite[Theorem~4.5]{samuelquantum} and gave the corresponding
algorithms \cite[Methods~4.9--4.11]{samuelquantum}.
We also established the quantum and doubled Pieri formulas and the
quantum expansions \cite[Theorems~4.6--4.7]{samuelquantum}, which we use
in Section~\ref{sec:quantum}.
The contribution of the present paper is an explicit accounting of the
full-product algorithms in parameters obtainable from their inputs.

Pak and Slonim \cite[Theorem~1.4]{pakslonim} prove that a
\emph{single ordinary} coefficient $c_{u,v}^w$, for three permutations in
$S_n$ supplied by their range-restricted codes, can be computed in time
\begin{equation}\label{eq:pak}
 \exp\bigl(O(\psi^{5/2})\bigr)
       (\log n)^{O(\psi^2\log\psi)}.
\end{equation}
Here $\psi$ is their total-range parameter.
The bound is polynomial in $\log n$ for fixed $\psi$.
Their method expresses a coefficient as an alternating sum of
lattice-point counts and applies Barvinok's algorithm in bounded
dimension. The same polyhedral construction proves eventual
quasi-polynomiality of stretched Schubert coefficients
\cite[Theorem~1.1]{pakslonim}.

Our parameter $r$ measures a different feature: how far the nonzero
Lehmer code extends from its first position.
It must not be identified with the number of nonzero entries, the number
of descents, or a translated window of width $\psi$.
We quote \eqref{eq:pak} with the parameter conventions of its source,
rather than deriving a formula relating $\psi$ to $r$.\footnote{The
version of \cite{pakslonim} consulted defines its range using exceedances
$w(i)>i$ and also states that the code vanishes off the exceedance set.
These two supports are not generally identical: $321$ has code $(2,1,0)$
and exceedance set $\{1\}$.
Our parameters and proofs use the Lehmer code directly and do not use
that identification.}

There are three important differences between \eqref{eq:pak} and
Theorem~\ref{thm:main}.
The former returns one integer.
The latter returns every coefficient of the triple product, each a
polynomial in two alphabets, within the same bound that it gives for
the ordinary integer coefficients.
The former is a bit-complexity statement for compressed input; the latter
initially counts arithmetic and elementary operations.
Finally, polynomial dependence on $n$ is much weaker than polynomial
dependence on $\log n$.
Thus the new bound does not subsume the bounded-range theorem.
At the unrestricted extreme our $2^{O(n\log n)}$ bound is smaller in
order than the bound obtained by substituting $\psi$ of order $n$
in \eqref{eq:pak}; this comparison of upper bounds is not a lower bound
on their algorithm or a claim about practical performance.

Encoding also matters for hardness statements.
Narayanan proves that Littlewood--Richardson coefficient computation is
$\#P$-complete for partitions encoded in binary
\cite{narayanan2006complexity}.
The associated Grassmannian Lehmer codes have a polynomial-size binary
encoding, so this gives $\#P$-hardness for the correspondingly compressed
Schubert problem.
It does not automatically give hardness for one-line permutations:
writing down a permutation associated with a large binary partition
part can require exponentially many entries.
Pak and Robichaux discuss this distinction and the unresolved
unary-input hardness question \cite{pak2024vanishing}.
Their subsequent result places Schubert vanishing in probabilistic
polynomial time \cite{pakrobichaux2025bpp}, improving the earlier
polynomial-hierarchy result.
Neither vanishing nor a single-coefficient hardness statement is used
in our running-time proof.

\section{Dominant completion and a justified label filter}\label{sec:completion}

A permutation is \emph{dominant} if its code is a partition.
Its Rothe diagram is then a Young diagram.
If $\theta=\code(\tau^{-1})$ is its list of column heights, the
dominant factorization is
\begin{equation}\label{eq:dominant}
 \Sch_\tau(x;t)=\prod_{i=1}^{L}\prod_{j=1}^{\theta_i}(x_j-t_i).
\end{equation}
In particular, the polynomial on the right is indexed by $\tau$,
not by $\tau^{-1}$. This distinction is essential.

\begin{method}[Completion]\label{method:completion}
Start with $\mu=v^{-1}$.
While its code has a strict adjacent ascent
$\code_i(\mu)<\code_{i+1}(\mu)$, choose the largest such $i$
and replace $\mu$ by $\mu s_i$.
At termination put $\theta=\code(\mu)$ with trailing zeros deleted,
$\tau=\mu^{-1}$, and $a_*=v\mu$.
\end{method}

\begin{lemma}\label{lem:completion}
The completion remains in $S_n$ and terminates at a dominant permutation.
It satisfies
\[
 \ell(a_*)=\ell(\tau)-\ell(v),\qquad
 \Des(\tau)=\Des(v),\qquad
 \theta_1=r_v,\qquad L\leq n-1
\]
when $v\ne\id$.
It can be constructed in $O(n^3)$ elementary operations.
\end{lemma}
\begin{proof}
For an adjacent code ascent, right multiplication by $s_i$ changes
the two code entries $(a,b)$ to $(b+1,a)$ and increases length by one.
Thus there are at most $\binom n2$ steps and the final code is weakly
decreasing.
The successive right multiplications give a length-additive
factorization $\mu=v^{-1}a_*$, proving the length identity.

To check the descent assertion, consider the inverses before and after
one step. They differ by interchanging the consecutive values $i,i+1$.
Their positions are $\mu(i),\mu(i+1)$, with $\mu(i)<\mu(i+1)$.
These positions cannot be adjacent: if they were, the two code entries
of $\mu$ would be equal. Swapping consecutive values in nonadjacent
positions does not change any adjacent comparison.
The descent set of the inverse therefore stays unchanged throughout.
For a dominant permutation the number of nonzero rows of the diagram
is its last descent; transposing the diagram identifies this number
with $\theta_1$.
Finally, the diagram of any permutation in $S_n$ is contained in the
staircase, so $L\leq n-1$.
Maintaining the code, scanning for the next ascent, and updating the
permutation costs $O(n)$ per step.
\end{proof}

For $0\leq i\leq L$ define the dominant permutation $\tau_i$ by
\[
 \code(\tau_i^{-1})=(\theta_1,\ldots,\theta_i),\qquad
 F_i(x;t)=\Sch_{\tau_i}(x;t).
\]
Here $\tau_0=\id$, and
$F_i=F_{i-1}\prod_{j=1}^{\theta_i}(x_j-t_i)$.
All these permutations are in $S_n$.
Use the signed double-Schubert basis in the auxiliary variables
\[
 T_a(t;z)=(-1)^{\ell(a)}\Sch_a(t;z).
\]
We say $\sigma\leq_L\tau$ if $\tau=\rho\sigma$ with
$\ell(\tau)=\ell(\rho)+\ell(\sigma)$.

\begin{lemma}[Cauchy filter]\label{lem:filter}
Set
\[
 A_i=\{a\in S_n:\ell(a\tau_i)=\ell(\tau_i)-\ell(a)\}.
\]
Then
\begin{equation}\label{eq:filter}
 F_i(x;t)=\sum_{a\in A_i}\Sch_{a\tau_i}(x;z)T_a(t;z).
\end{equation}
Consequently $A_0=\{\id\}$, $a_*\in A_L$, and the coefficient of
$T_{a_*}$ in $F_L$ is $\Sch_v(x;z)$.
Moreover,
\begin{equation}\label{eq:labelcount}
 |A_i|\leq \fall{n}{h}=V.
\end{equation}
\end{lemma}
\begin{proof}
The double Cauchy identity and the interchange-of-alphabets identity give
\[
 \Sch_{\tau_i}(x;t)
 =\sum_{\substack{\rho\sigma=\tau_i\\ \ell(\rho)+\ell(\sigma)=\ell(\tau_i)}}
       \Sch_\sigma(x;z)\Sch_\rho(z;t),\qquad
 \Sch_\rho(z;t)=(-1)^{\ell(\rho)}\Sch_{\rho^{-1}}(t;z).
\]
Set $a=\rho^{-1}$. Then $\sigma=a\tau_i$ and these are exactly
the terms in \eqref{eq:filter}.
The claims about $a_*$ follow from Lemma~\ref{lem:completion}.

For the count, the map $a\mapsto a\tau_i$ is a bijection from $A_i$
onto the left weak interval $[\id,\tau_i]_L$.
If $\sigma\leq_L\tau_i$, every right descent of $\sigma$ is a
right descent of $\tau_i$: writing $\tau_i=\rho\sigma$,
a descent of $\sigma$ gives
$\ell(\tau_i s_j)\leq\ell(\rho)+\ell(\sigma)-1$.
The diagram of $\tau_i$ has at most $h$ rows, so these permutations
have no descent after position $h$.
A permutation in $S_n$ with an increasing suffix after position $h$
is determined by its ordered first $h$ values.
There are exactly $\fall{n}{h}$ such permutations.
\end{proof}

The filter is a polynomial identity, not an assertion that each discarded
path has zero weight. Coefficients of labels outside $A_i$ can vanish by
cancellation. They may be discarded \emph{at a column boundary} because
their entire coefficient is zero and subsequent multiplication is linear.
In particular, this does not justify pruning the interior of a column
merely because an intermediate path node fails the boundary test.

\section{A column expansion with a direct proof}\label{sec:columns}

For an ordered alphabet $Y=(Y_1,\ldots,Y_m)$ define
\begin{equation}\label{eq:elementary}
 E_{p,m}(Y;Z)=
 \sum_{1\leq j_1<\cdots<j_p\leq m}
       \prod_{q=1}^p (Y_{j_q}-Z_{j_q-q+1}),
 \qquad 0\leq p\leq m .
\end{equation}
Only the first $m-p+1$ entries of $Z$ are used when $p>0$.
Set $E_{0,m}=1$ and $E_{p,m}=0$ for $p>m$.
Repeated entries in $Z$ are allowed.

A \emph{pivot-$i$ path} of length $b$ is a sequence
$P=(a_0,\ldots,a_b)$ with
\[
 a_j=a_{j-1}t_{i,c_j},\qquad c_j\ne i,\qquad
 \ell(a_j)=\ell(a_{j-1})+1.
\]
No distinctness condition is imposed on the $c_j$.
Associate to it
\[
 \varepsilon(P)=(-1)^{\#\{j:c_j<i\}},\qquad
 Z(P)=(z_{a_0(i)},z_{a_1(i)},\ldots,z_{a_b(i)}).
\]

\begin{lemma}[Unrestricted column formula]\label{lem:column}
For $C_i=\prod_{j=1}^k(x_j-t_i)$,
\begin{equation}\label{eq:column}
 C_i T_a(t;z)=
 \sum_{b=0}^k\ \sum_{P:a_0=a}
   \varepsilon(P)E_{k-b,k}(x_1,\ldots,x_k;Z(P))T_{a_b}(t;z).
\end{equation}
\end{lemma}
\begin{proof}
The equivariant single-variable Monk formula, in the signed basis, is
\begin{equation}\label{eq:monk}
 (X-t_i)T_a=(X-z_{a(i)})T_a
   +\sum_{\substack{j>i\\ \ell(at_{ij})=\ell(a)+1}}T_{at_{ij}}
   -\sum_{\substack{j<i\\ \ell(at_{ji})=\ell(a)+1}}T_{at_{ji}} .
\end{equation}
It follows by subtracting the formulas for multiplication by
$t_1+\cdots+t_i$ and $t_1+\cdots+t_{i-1}$.
Apply \eqref{eq:monk} successively with $X=x_1,\ldots,x_k$.
Fix the subsequence of non-diagonal steps, namely a path $P$ of length $b$.
Its sign is $\varepsilon(P)$.
Choose the $k-b$ diagonal positions
$j_1<\cdots<j_{k-b}$ among the $k$ factors.
Before the diagonal position $j_q$ there have been $j_q-q$
non-diagonal steps, so its contribution is
$x_{j_q}-z_{a_{j_q-q}(i)}$.
Summing over the diagonal positions gives exactly
$E_{k-b,k}(x;Z(P))$ in \eqref{eq:elementary}.
\end{proof}

\begin{lemma}[Distinct-position cancellation]\label{lem:cancellation}
The sum in \eqref{eq:column} is unchanged if it is restricted to paths
whose nonpivot positions $c_1,\ldots,c_b$ are pairwise distinct.
For these paths $Z(P)$ consists of the values
\[
 D_i(a_0,a_b)=\{a_0(i)\}\cup
              \{a_0(j):a_0(j)\ne a_b(j)\},
\]
each once, and its order is immaterial.
\end{lemma}
\begin{proof}
This is the cancellation underlying
\cite[Lemma~7.5 and Theorem~7.6]{samuelquantum}.
We explain how it applies to the signed-basis convention used here.
In a path with a repetition, take the segment ending at the first
repeated nonpivot position, starting at its earlier occurrence.
The rotation lemma pairs its valid cyclic rotations with starting
positions on opposite sides of the pivot.
The resulting involution fixes the two endpoints, the path length,
and the multiset of values encountered at the pivot, and reverses
$\varepsilon(P)$.
By \cite[Lemma~7.3]{samuelquantum},
$E_{k-b,k}(x;Z)$ is symmetric in the $b+1$ entries of $Z$.
The two weights therefore cancel.
The sign convention in the cited skew-divided-difference formula
differs by $(-1)^b$, which is constant on each such pair.

On a distinct-position path the pivot successively carries
$a_0(i),a_0(c_1),\ldots,a_0(c_b)$.
These are exactly the values in $D_i(a_0,a_b)$.
The same symmetry gives the last assertion.
\end{proof}

\begin{theorem}[Filtered elementary expansion]\label{thm:elementary}
Start at label $\id$ and, for $i=1,\ldots,L$, take a pivot-$i$
path $P_i$ of length $b_i\leq\theta_i$ with pairwise distinct
nonpivot positions within that path.
Its starting label is the previous path's endpoint.
Require the endpoint at column $i$ to belong to $A_i$ and the final
endpoint to be $a_*$.
Then
\begin{equation}\label{eq:expansion}
 \Sch_v(x;z)=
 \sum_{(P_1,\ldots,P_L)}
 \prod_{i=1}^L
   \varepsilon(P_i)E_{\theta_i-b_i,\theta_i}(x;Z(P_i)).
\end{equation}
\end{theorem}
\begin{proof}
Apply Lemmas~\ref{lem:column} and \ref{lem:cancellation}
to the factors of $F_L$.
At each column boundary, Lemma~\ref{lem:filter} permits deleting all
labels outside $A_i$.
Extracting $T_{a_*}$ gives \eqref{eq:expansion}.
A path reaching $a_*\in S_n$ never leaves $S_n$, since all its
steps increase Bruhat order and the interval below a permutation in
$S_n$ is contained in $S_n$.
Before imposing the boundary filters, this is precisely
\cite[Theorem~4.5]{samuelquantum}.
\end{proof}

Setting $z=0$ in \eqref{eq:expansion} replaces $E_{p,m}$ by
$e_p(x_1,\ldots,x_m)$ and preserves the paths and their signs.
This is a statement about the \emph{path-indexed} expansion.
It is not a way of recovering the double polynomial from a collected
ordinary expansion after its path labels and alphabets have been erased.
No uniqueness or minimality of an elementary representation is needed.

\begin{method}[Backward construction of the label graphs]\label{method:labels}
Start with the singleton set $\{a_*\}$ at boundary $L$.
For $i=L,L-1,\ldots,1$ and each currently retained upper label,
enumerate all descending pivot-$i$ Bruhat paths of lengths
$0,\ldots,\theta_i$, never reusing a nonpivot position in a path.
Retain a path if its lower endpoint is in $A_{i-1}$.
Reverse it to obtain a forward edge, storing its two endpoints,
length, sign, and ordered alphabet indices.
Collect the distinct lower endpoints to seed the next boundary.
Do not prune internal path nodes using the endpoint filter.
\end{method}

This construction gives exactly the paths in Theorem~\ref{thm:elementary}
that can reach the required final label; it does not enumerate every
permutation in $A_i$.
Each stored edge represents an entire column path, not a single Bruhat cover.
The product of distinct pivot transpositions is a cycle, from which
their order can be recovered, so at most one such path has a specified
pair of endpoints.
This is the \texttt{kdown}/\texttt{vpaths} construction of
\cite[Methods~4.9--4.10]{samuelquantum}.
In its implementation, a nonpivot position is eligible exactly when it
still holds its value from the start of the path: distinct pivot swaps
change every position already used and leave all others untouched.

\section{Pieri multiplication and coefficient construction}\label{sec:pieri}

For a permutation $w$ and an integer $k$, a canonical elementary
Pieri chain consists of Bruhat covers
$w_j=w_{j-1}t_{a_j b_j}$ with $a_j\leq k<b_j$,
whose moved-in values $w_{j-1}(b_j)$ are strictly decreasing.
The elementary Pieri rule identifies chains of length $a$ with the
terms of $e_a(x_1,\ldots,x_k)\Sch_w(x)$, each with multiplicity one
\cite{sottile}.
In particular, distinct canonical chains from a fixed $w$ have distinct
endpoints, including across different lengths.
No left position can be used twice: a second ascent there would require
a moved-in value greater than the first, contrary to the ordering.

Write $w\longrightarrow_k w'$ for such an endpoint and let
$a=\ell(w')-\ell(w)$.
The fixed alphabet is
\[
 Y(w,w';k)=(y_{w(j)}:1\leq j\leq k,\ w(j)=w'(j)),
\]
listed in increasing order of $j$; it has $k-a$ entries.
The double Pieri formula \cite[Theorem~7.1]{samuelmolev} is
\begin{equation}\label{eq:pieri}
 \Sch_w(x;y)E_{p,k}(x;Z)
 =\sum_{\substack{w\longrightarrow_k w'\\ a\leq p}}
     E_{p-a,k-a}(Y(w,w';k);Z)\Sch_{w'}(x;y).
\end{equation}
Since this is a polynomial identity in the coefficient alphabet,
it remains valid when entries of $Z$ coincide, as they may in
Lemma~\ref{lem:column}.

\begin{method}[Pieri enumeration]\label{method:pieri}
Maintain a frontier $(w,\beta)$, initially $(w,+\infty)$.
To extend it, test $a\leq k<b$ for
\[
 w(a)<w(b)<\beta,\qquad
 \nexists c\ (a<c<b,\ w(a)<w(c)<w(b)).
\]
For each successful test append $(wt_{ab},w(b))$ to the next frontier.
Return the initial permutation and every endpoint through the requested
maximum degree $p\leq k$, recording its degree.
\end{method}

\begin{proposition}\label{prop:pieriwork}
If all generated permutations lie in $S_M$ and $B$ endpoints are
returned, Method~\ref{method:pieri} costs $O(M^3B)$ elementary operations.
\end{proposition}
\begin{proof}
A frontier node has at most $k(M-k)$ candidate transpositions.
Testing the Bruhat-cover condition and copying a permutation each cost
$O(M)$.
Multiplicity-freeness of the canonical Pieri chains bounds the total
number of processed nodes by the number of returned endpoints.
\end{proof}

\begin{lemma}[Polynomial-cost elementary weights]\label{lem:weight}
A factor $E_{p,m}(Y;Z)$ can be built as an arithmetic circuit using
$O((m+1)^2)$ ring operations.
\end{lemma}
\begin{proof}
Separating subsets according to whether they contain the last index gives
\begin{equation}\label{eq:weightrec}
 E_{p,m}(Y;Z)=E_{p,m-1}(Y_1,\ldots,Y_{m-1};Z)
 +(Y_m-Z_{m-p+1})
     E_{p-1,m-1}(Y_1,\ldots,Y_{m-1};Z).
\end{equation}
Use $E_{0,t}=1$ and $E_{q,t}=0$ for $q>t$, and memoize the
subproblems $(q,t)$ reached from $(p,m)$.
There are at most $(m+1)^2$ of them, each costing a bounded number
of operations.
Every nontrivial reached subproblem has $t-q\leq m-p$,
so no coefficient variable beyond $Z_{m-p+1}$ is requested.
The circuit shares subexpressions rather than copying expression trees.
\end{proof}

An unmemoized divide-and-conquer evaluation is another valid choice:
split $Y$ after $s=\lfloor m/2\rfloor$ and use
\[
 E_{p,m}(Y;Z)=\sum_a E_{a,s}(Y';Z)
     E_{p-a,m-s}(Y'';Z^{[s-a]}),\qquad
 Z^{[c]}_j=Z_{j+c},
\]
where the sum is over feasible degrees.
Its recursion gives an upper bound $m^{O(\log m)}$.
That is an evaluation-strategy cost, not an intrinsic lower bound
for factorial elementary polynomials.
Theorem~\ref{thm:main} uses \eqref{eq:weightrec}.

\begin{method}[Classical \alg]\label{method:product}
Construct the completion and the label graphs.
The coefficient ring is $\mathbb Z[y,z]$.
Initialize a table $H$ with $H(u,\id)=1$.
At column $i$ put $k=\theta_i$ and create a fresh table $H'$.
\begin{enumerate}
\item Group the keys of $H$ by their first component $w$.
For each distinct $w$ enumerate its Pieri endpoints once, through degree
$\min(k,d-\ell(w))$.
Reuse this list for all labels paired with $w$.
\item For each table entry $H(w,a)$, each label edge $P:a\to a'$
of length $b$, and each Pieri endpoint $w'$ of degree $a_u$,
use the pair only if $a_u+b\leq k$.
Accumulate
\[
 H'(w',a')\mathrel{+}=
  H(w,a)\varepsilon(P)
  E_{k-a_u-b,k-a_u}\bigl(Y(w,w';k);Z(P)\bigr).
\]
\item Replace $H$ by $H'$ and continue.
After column $L$, return $H(w,a_*)$ as the coefficient of $\Sch_w(x;y)$.
\end{enumerate}
In the ordinary specialization keep only $a_u+b=k$;
the elementary weight is then $1$.
\end{method}

\begin{theorem}[Correctness]\label{thm:correct}
Method~\ref{method:product} computes \eqref{eq:product}.
\end{theorem}
\begin{proof}
Multiply the expansion in Theorem~\ref{thm:elementary} by
$\Sch_u(x;y)$ and apply \eqref{eq:pieri} at each column.
The displayed update is exactly the product of the column sign,
the running coefficient, and the Pieri coefficient.
Dynamic programming merely combines contributions with equal endpoint
pairs.
The final label condition is precisely that in
Theorem~\ref{thm:elementary}.

It remains to justify the degree cap.
Along a complete label path,
$\sum_i b_i=\ell(a_*)=\ell(\tau)-\ell(v)$ and
$\sum_i\theta_i=\ell(\tau)$.
Every admissible paired path has
$\sum_i a_{u,i}\leq\sum_i(\theta_i-b_i)=\ell(v)$.
All $a_{u,i}$ are nonnegative, so no contributing prefix can reach
a permutation of length greater than $d$.
Thus the cap discards no contribution.
Specialization to $y=z=0$ annihilates every elementary weight of
positive degree and leaves precisely the stated ordinary update.
\end{proof}

\section{Input bounds for every enumerated state}\label{sec:states}

\begin{lemma}[Pieri ambient size]\label{lem:ambient}
If $w\in S_m$ and $k<m$, a canonical elementary Pieri chain from
$w$ stays in $S_{m+1}$.
\end{lemma}
\begin{proof}
Before any position beyond $m$ is used, a Bruhat cover cannot have
right endpoint $b\geq m+2$: position $m+1$ has value $m+1$,
strictly between $w(a)\leq m$ and $w(b)=b$.
The only possible new value is therefore $m+1$.
If it is moved in, it must be the first moved-in value, since all
previously available values are at most $m$ and the sequence decreases.
All subsequent moved-in values are smaller; no larger new position
can be used.
\end{proof}

\begin{lemma}[Code-length and degree bounds]\label{lem:states}
Every $u$-side permutation enumerated by Method~\ref{method:product},
including nodes of Pieri searches that do not survive into the next
table, lies in $S_M$, has no descent after $r$, and has length at most $d$.
There are at most $U$ such permutations.
\end{lemma}
\begin{proof}
The ambient claim follows by applying Lemma~\ref{lem:ambient} once
per column, starting in $S_n$.
The degree cap in Method~\ref{method:product} bounds every partial
Pieri chain by $d$.

For the descent claim, let $w$ have no descent after $r$.
Since $k\leq h\leq r$, the polynomial
\[
 e_a(x_1,\ldots,x_k)\Sch_w(x)
\]
involves only $x_1,\ldots,x_r$.
For $j>r$ its divided difference $\partial_j$ is zero.
For a Schubert basis element, $\partial_j\Sch_b=\Sch_{bs_j}$
if $j\in\Des(b)$; otherwise its divided difference is zero.
Distinct $b$ with descent $j$ give distinct $bs_j$.
Linear independence of the Schubert basis therefore implies that no
permutation in this product has a descent after $r$.
Every canonical chain of length $a$ is a term of this ordinary Pieri
product with coefficient one, so this argument applies to
every enumerated endpoint, before any double-coefficient cancellation.
Induction starts with $u$.

Permutations of $S_M$ with an increasing suffix after $r$ number
$\fall{M}{r}$.
Alternatively, their Lehmer codes have length at most $r$ and sum
at most $d$. There are $\binom{d+r}{r}$ nonnegative $r$-tuples with
that sum bound, by stars and bars.
Taking the smaller estimate proves the assertion.
\end{proof}

\begin{lemma}[Cost of the label construction]\label{lem:labelwork}
Method~\ref{method:labels} stores at most $L V Q_h(n)$ column paths
and costs $O(n^3+L n^3 V Q_h(n))$ elementary operations.
This estimate charges paths rejected by the endpoint filter.
\end{lemma}
\begin{proof}
There are at most $V$ retained upper labels at a boundary, by
Lemma~\ref{lem:filter}.
From each upper label, a descending pivot path has depth at most
$\theta_i\leq h$ and uses each nonpivot position at most once.
There are at most $\fall{n-1}{b}$ possible position sequences of
length $b$, even before testing the Bruhat conditions.
The entire search tree, including rejected nodes, consequently has
at most $Q_h(n)$ nodes.
Every node is in $S_n$.
Scanning candidates, testing Bruhat covers, computing the endpoint
length test, and recording a path take polynomial work bounded by
$O(n^3)$ per node.
Summing over boundaries and upper labels proves the claim.
Constructing the prefixes $\tau_i$ and the completion fits the
additional $O(n^3)$ term.
\end{proof}

These estimates avoid two potentially misleading shortcuts:
counting only retained endpoints does not bound the work of a search
that explores rejected paths, and controlling final support does not
by itself control intermediate states.
Both kinds of work have been bounded explicitly.

\section{Complexity and its interpretation}\label{sec:complexity}

\subsection{Arithmetic model}

A coefficient is stored as a directed acyclic arithmetic circuit over
$\mathbb Z[y,z]$; addition and multiplication create gates pointing to
existing gates.
The elementary routine shares its dynamic-programming subexpressions.
We count these ring operations and elementary operations on permutation
entries separately from bit costs.
Permutation entries have $O(\log M)$ bits.
Deterministic dictionaries may be implemented by balanced trees:
comparing two keys costs $O(M)$ elementary operations, and a table
has at most $UV$ keys.
Since $\log(UV)=O(M\log M)$, this bookkeeping is absorbed by the
$M^3$ factor below.
No unit-cost polynomial identity test is assumed.

The returned table is an exact expansion but need not be a
minimal-support list.
An identically zero coefficient can be left as a circuit computing zero.
Removing all such entries, expanding coefficients into monomials, and
canonical symbolic simplification are different output requirements
and are not included in the theorem.

\begin{proof}[Proof of Theorem~\ref{thm:main}]
The label preprocessing costs
$O(n^3+L n^3 VQ_h(n))$ by Lemma~\ref{lem:labelwork}.
At a product boundary there are at most $U$ distinct $u$-side
permutations and at most $V$ labels.
At one column the Pieri enumeration is called at most $U$ times,
each returning at most $U$ endpoints by Lemma~\ref{lem:states}.
Proposition~\ref{prop:pieriwork} bounds this work by $O(M^3U^2)$.

Count matched pairs by iterating over all column edges, rather than
assuming a bound on the forward outdegree of each label.
There are at most $VQ_h(n)$ such edges in a column.
For each there are at most $U$ possible source permutations,
and each source has at most $U$ Pieri endpoints.
Thus there are at most $U^2VQ_h(n)$ pairs.
Forming the fixed alphabet takes $O(M)$ work.
Lemma~\ref{lem:weight} builds its elementary coefficient in
$O((h+1)^2)$ ring operations; the remaining update takes a bounded
number of ring operations.
Permutation and dictionary work, as well as this arithmetic, is
bounded by $O(M^3)$ per pair.
Summing over $L$ columns and absorbing the preprocessing and
Pieri terms gives \eqref{eq:main}.

For the coarse estimate use
\[
 L\leq n-1,\quad M\leq2n-1,\quad U\leq(2n)^r,\quad
 V\leq n^h,\quad Q_h(n)\leq2n^h .
 \]
The last bound follows by comparing $\fall{n-1}{b}$ with $(n-1)^b$
and summing the geometric series for $n\geq3$, and
for $n=2$ because $h\leq1$.
Substitution gives \eqref{eq:coarse}.
Since $h\leq r\leq n-1$, its logarithm is at most
$4n\log_2 n+O(n)$.

If $v$ is dominant, $\mu=v^{-1}$ and $a_*=\id$.
No descending path of positive length starts at the identity.
Every label graph therefore consists of the single degree-zero edge
$\id\to\id$, and the entire product calculation has one label.
Replacing $VQ_h(n)$ by $1$ gives the dominant-input estimate.
By Proposition~\ref{prop:uniform}, the double and ordinary products
satisfy the same bound.
\end{proof}

\subsection{Uniformity across coefficient rings}

The ordinary specialization of Method~\ref{method:product} keeps only
pairs with $a_u+b=k$. In the double specialization, set $z=y$.

\begin{proposition}[Uniformity]\label{prop:uniform}
Fix $u,v$ with $v\ne\id$.
Compare the triple computation of Method~\ref{method:product} with its
double and ordinary specializations.
\begin{enumerate}
\item All three use the same completion and the same label graphs.
\item Every table key and every Pieri call made by the ordinary
computation is made by the triple computation, with the same arguments.
Every matched pair used by the ordinary computation is used by the
triple computation.
The triple computation's additional pairs are those with $a_u+b<k$.
\item Per matched pair, the triple and double computations perform
$O((h+1)^2)$ ring operations and $O(M)$ elementary operations beyond
the ordinary computation.
\end{enumerate}
In particular, all three computations satisfy \eqref{eq:main}.
The same statements hold for Method~\ref{method:qproduct}, with
$t+b=k$ in place of $a_u+b=k$ and $\tilde h$ in place of $h$, except
that a $\mathtt{qpieri}$ call of the triple computation may have a
larger depth $p_w$. Its returned list then contains that of the
ordinary call.
\end{proposition}
\begin{proof}
The completion and the label graphs depend only on $v$.
The recorded alphabet indices $Z(P)$ are used only to form weights,
never to accept or reject a path.
Step~(1) of Method~\ref{method:product} depends only on $w$, $k$, and $d$.
Step~(2) accepts a pair according to $a_u$ and $b$ alone.
By induction over columns, the ordinary keys at each boundary therefore
form a subset of the triple keys, and the ordinary pairs form a subset
of the triple pairs.
The additional work for a triple pair is the formation of $Y(w,w';k)$
and the circuit of Lemma~\ref{lem:weight}.
In the ordinary computation the weight is $1$.
The proofs of Lemmas~\ref{lem:filter}, \ref{lem:states}, and
\ref{lem:labelwork} never use the coefficient ring, so they apply to all
three computations.
For the quantum method, the same argument applies with
Lemmas~\ref{lem:strict} and \ref{lem:qstates}.
A $\mathtt{qpieri}$ search of depth $p$ is a prefix of one of larger
depth, and Lemma~\ref{lem:qstates} bounds every chain that the
triple computation enumerates.
\end{proof}

The proposition is an upper-bound statement.
The ordinary computation can be cheaper in practice, since it discards
pairs with $a_u+b<k$.
The point is that the worst-case analysis does not distinguish the
three problems, and the bound for the triple product is not obtained by
paying for its extra generality.

\begin{corollary}[Ordinary bit costs]\label{cor:bits}
For unit initial coefficient, the ordinary algorithm has
bit complexity bounded by \eqref{eq:main} times a polynomial in $n$.
It therefore still has complexity $2^{O(n\log n)}$, and polynomial
complexity in $n$ for fixed $r_u,r_v$.
\end{corollary}
\begin{proof}
Each surviving matched pair contributes the running integer with
sign $+1$ or $-1$.
The absolute value of any intermediate sum is bounded by the number
of complete or partial paired paths.
A state has at most $UVQ_h(n)$ successors counted with multiplicity,
so that number is at most $(2UVQ_h(n))^L$.
Its bit length is
\[
 O\bigl(1+L\log(2UVQ_h(n))\bigr)=O(n^2\log n).
 \]
Integer additions and permutation operations consequently incur only
a polynomial bit overhead.
This argument bounds intermediate signed sums as well as final
nonnegative coefficients; it does not rely on positivity preventing
intermediate cancellation.
\end{proof}

\subsection{What bounded code length does and does not imply}

The number of code entries is a useful parameter because the
$u$-side count can be replaced by $\binom{d+r}{r}$ when degree
is more restrictive than the ambient permutation size.
No weak-order interval enumeration is needed for this refinement.
The $v$-side bound depends on $h=r_v$, reflecting the fact that the
second factor supplies the elementary expansion.
For ordinary multiplication the roles may be interchanged and the
smaller of the two bounds used.
For \eqref{eq:product}, such an interchange also changes which
coefficient alphabet indexes the output basis and is not automatic.

An input code with $r_w$ entries has encoding length
$O(r_w\log n)$, whereas the algorithm explicitly works with
permutations and has polynomial dependence on $n$ for fixed $r$.
This is a pseudo-polynomial feature relative to binary code entries.
For example, $\Sch_{\operatorname{uncode}(0,A)}(x_1,x_2)=h_A(x_1,x_2)$,
and
\[
 h_A(x_1,x_2)^2
   =\sum_{j=0}^{A}s_{(2A-j,j)}(x_1,x_2).
\]
The two input codes have constant length and $O(\log(A+1))$ bits,
but an explicit Schubert-basis list has $A+1$ terms.
This example explains why a full-product algorithm and a single-coefficient
algorithm have fundamentally different output obligations; it is not
used as a running-time parameter.

Under scaling of the entries of fixed Lehmer codes by a factor $N$,
the code lengths remain fixed and $n=O(N)$.
Theorem~\ref{thm:main} hence gives a polynomial-in-$N$ bound for these
full products.
This does not yield the polylogarithmic dependence available from
\eqref{eq:pak} in its fixed-range regime, nor does it prove eventual
quasi-polynomiality.
Finite computations or interpolation without a proved threshold and
period do not establish that conclusion.

\begin{corollary}[Grassmannian inputs]\label{cor:grass}
Let $u,v\in S_n$, $v\ne\id$, have no descents other than at $k$.
Then every column of the completion has height $k$, and every $u$-side
permutation enumerated by Method~\ref{method:product}, including
Pieri search nodes, is $k$-Grassmannian.
Every label $a\in A_i$ has $a\tau_i$ $k$-Grassmannian.
Consequently $U$ and $V$ in \eqref{eq:main} may be replaced by
\[
 U_k=\min\left\{\binom{M}{k},\binom{d+k}{k}\right\},\qquad
 V_k=\binom nk ,
\]
giving
\[
 O\left(n^3+LM^3U_k^2V_kQ_k(n)\right)
 =O\left(n^{4}\binom{2n}{k}^{2}\binom nk n^{k}\right).
\]
Uniformly in $k$ this is $O(4^kn^{4k+4}/(k!)^3)$, and in general it is
$2^{O(n)}n^{k+4}$.
\end{corollary}
\begin{proof}
By Lemma~\ref{lem:completion}, $\Des(\tau)=\Des(v)=\{k\}$.
A dominant permutation whose only descent is at $k$ has a rectangular
diagram with $k$ rows, so its transpose $\theta$ has all parts equal to
$k$.
For the labels, $a\tau_i\leq_L\tau_i$, so as in the proof of
Lemma~\ref{lem:filter} its descents lie in $\Des(\tau_i)=\{k\}$;
there are at most $\binom nk$ such permutations in $S_n$.

Let $w$ be $k$-Grassmannian. Then $\Sch_w(x)$ is symmetric in
$x_1,\ldots,x_k$ and involves no other $x$-variables, so the same
holds for $e_a(x_1,\ldots,x_k)\Sch_w(x)$.
Hence $\partial_j$ annihilates it for every $j\ne k$.
The argument of Lemma~\ref{lem:states} shows that every permutation in
this ordinary Pieri product, and so every canonical chain endpoint,
has no descent other than at $k$.
Induction from $u$ covers all enumerated states, including those of
the double computation, whose Pieri chains are the same.
A $k$-Grassmannian permutation in $S_M$ is determined by its first $k$
values, as a set, giving $\binom Mk$.
Alternatively, its code is a partition with at most $k$ parts,
read in reverse, and of size at most $d$; there are at most
$\binom{d+k}{k}$ of these.
The displayed estimate follows from the proof of Theorem~\ref{thm:main},
with $M\leq 2n$ and $Q_k(n)\leq 2n^k$.
\end{proof}

Thus Grassmannian multiplication works entirely with partitions in a
$k\times(M-k)$ box on the $u$-side, with
$L=\code_k(v)\leq n-k$ columns.
The remaining non-exponential factor $n^{k}$ comes from the label
search $Q_k(n)$, which counts rejected distinct-position paths.
Sharpening that count would make the Grassmannian bound $2^{O(n)}$.

\section{Quantum multiplication}\label{sec:quantum}

We now prove Theorem~\ref{thm:quantum}, directly for the triple quantum
product.
The architecture is unchanged: a filtered elementary expansion of the
second factor is paired with a Pieri rule for the first.
Two ingredients change.
The factors must be quantized one column at a time, which forces
pairwise distinct column heights; and the quantum Pieri rule has
length-decreasing steps, which requires new state estimates.

\subsection{Quantization and the strict completion}

Let $X_j=(x_1,\ldots,x_j)$.
A \emph{standard elementary monomial} is a product
$\prod_j e_{c_j}(X_j)$ in which each height $j$ occurs at most once and
$0\leq c_j\leq j$.
These monomials form a $\mathbb Z$-basis of $\mathbb Z[x]$.
The quantization map of Fomin, Gelfand, and Postnikov
\cite{fomin1997quantum} is the $\mathbb Z[q]$-linear map, extended
linearly over all coefficient alphabets, sending
$\prod_j e_{c_j}(X_j)$ to $\prod_j e^q_{c_j}(X_j)$, where
$e^q_c(X_j)$ are the quantum elementary polynomials.
It sends $\Sch_w(x;y)$ to the quantum double Schubert polynomial
$\widetilde\Sch_w(x;y)$ \cite{fomin1997quantum,kirillov2000quantum},
\cite[Lemma~2.1]{samuelquantum}.
These polynomials, indexed by $S_\infty$, form a
$\mathbb Z[q][y]$-basis of $\mathbb Z[q][x,y]$, so we may write
\begin{equation}\label{eq:qproduct}
 \widetilde\Sch_u(x;y)\widetilde\Sch_v(x;z)
  =\sum_{w,\delta}q^\delta c_{u,v}^{w,\delta}(y;z)\,
     \widetilde\Sch_w(x;y).
\end{equation}
At $z=y$ the coefficient $c_{u,v}^{w,\delta}(y;y)$ is the
equivariant quantum Littlewood--Richardson coefficient of $Fl_N$ for every
$N$ such that every occurring $w$ lies in $S_N$ and every occurring degree
has $\delta_i=0$ for $i\geq N$
\cite{kim1996equivariant,lam2014quantum,mihalcea2006equivariant},
\cite[Section~2]{samuelquantum}.
At $y=z=0$ it is the ordinary quantum structure constant.
Grading $\deg q_i=2$ makes \eqref{eq:qproduct} homogeneous.
Put
\[
 \widetilde E_{p,k}(x;Z)=\mathcal Q\bigl(E_{p,k}(x;Z)\bigr),
\]
where $\mathcal Q$ is the quantization map.
Expanding the definition, $E_{p,k}(x;Z)$ is a $\mathbb Z[Z]$-linear
combination of $e_c(X_k)$ with $c\leq k$, so $\widetilde E_{p,k}$ is
the corresponding combination of $e^q_c(X_k)$
\cite[Lemma~7.3]{samuelquantum}.

The classical completion of Section~\ref{sec:completion} usually has
repeated column heights, and $\mathcal Q$ is not multiplicative on
$e_a(X_k)e_b(X_k)$.
We therefore use the strict completion of
\cite[Definition~4.1]{samuelquantum}.
Call $\mu$ \emph{strict dominant} if it is dominant and its nonzero
code entries are strictly decreasing.

\begin{method}[Strict completion]\label{method:strict}
Start with $\mu=v^{-1}$.
While $\mu$ is not strict dominant, choose the largest $i$ with
$\code_{i+1}(\mu)\neq0$ and $\code_i(\mu)\leq\code_{i+1}(\mu)$,
and replace $\mu$ by $\mu s_i$.
At termination put $\tilde\theta=\code(\mu)$ with trailing zeros
deleted, $\tilde\tau=\mu^{-1}$, $\tilde a_*=v\mu$, and
$\tilde h=\tilde\theta_1$.
\end{method}

\begin{lemma}[Strict height]\label{lem:strict}
Let $v\ne\id$ and let $L$ be the number of nonzero entries of
$\code(v^{-1})$.
The strict completion stays in $S_n$, uses $O(n^3)$ elementary
operations, and satisfies
$\ell(\tilde a_*)=\ell(\tilde\tau)-\ell(v)$.
Its shape has exactly $L$ parts,
$\tilde\theta_1>\cdots>\tilde\theta_L>0$,
and the classical completion also has exactly $L$ columns.
Finally,
\begin{equation}\label{eq:strictheight}
 \max(r_v,L)\leq\tilde h\leq
 \min\bigl(n-1,\ r_v+L-1\bigr).
\end{equation}
\end{lemma}
\begin{proof}
A step at $i$ has $\code_i(\mu)\leq\code_{i+1}(\mu)$, so
$\mu(i)<\mu(i+1)$; it replaces the entries $(a,b)$ by $(b+1,a)$ and
increases length by one.
Since $b\neq0$, we have $i+1\leq m-1$ whenever $\mu\in S_m$, so
$\mu s_i\in S_m$.
This proves the ambient claim for $m=n$, termination, the length identity,
and the $O(n^3)$ cost as in Lemma~\ref{lem:completion}.
The replacement $(a,b)\mapsto(b+1,a)$ with $b\neq0$ preserves the
number of nonzero entries. The same holds for the classical steps,
which have $a<b$.
Hence both shapes have $L$ parts, and the termination condition gives
strictly decreasing parts.

A strictly decreasing sequence of $L$ positive integers has first entry at
least $L$.
For the other lower bound, $\tilde\tau=x^{-1}v$ is length-additive,
so, as in the proof of Lemma~\ref{lem:filter}, every right descent of
$v$ is one of $\tilde\tau$.
The permutation $\tilde\tau$ is the inverse of a dominant permutation,
hence dominant, and its last descent is the number of rows of its
diagram, namely $\tilde\theta_1$.
Thus $r_v\leq\tilde h$.

For the upper bound let $\mu$ be the classical completion, with shape
$\theta$, and put $m=\max_j(\theta_j+j)$.
A dominant permutation with shape $\theta$ lies in $S_m$:
its code satisfies $\code_j\leq m-j$.
Since $\mu=v^{-1}x$ with lengths adding, $v^{-1}$ is a prefix of a
reduced word of $\mu$ and so lies in $S_m$.
Applying the first paragraph with this $m$ shows that the strict
completion lies in $S_m$, whence $\tilde h=\code_1\leq m-1$.
Finally, $\theta_j\leq\theta_1=r_v$ and $j\leq L$ give
$m\leq r_v+L$. The bound $\tilde h\leq n-1$ is the case $m=n$.
\end{proof}

Both inequalities in \eqref{eq:strictheight} can be sharp.
For $v=(m+1)12\cdots m$, which has $\Sch_v=x_1^m$, one has $r_v=1$,
$L=m$, $\theta=(1^m)$, and $\tilde\theta=(m,m-1,\ldots,1)$.
In finite checks through $S_8$ we also observed
$\tilde h=\max_j(\theta_j+j)-1$ for every nonidentity $v$.
We do not use this equality.

\subsection{Quantum filtered expansion}

Define $\tilde\tau_i$, $\tilde F_i$, $\tilde A_i$, and the label graphs
exactly as in Sections~\ref{sec:completion} and \ref{sec:columns},
with $\tilde\theta$ in place of $\theta$, and put
\[
 \widetilde V=\fall{n}{\tilde h}.
\]
The proofs of Lemmas~\ref{lem:filter} and \ref{lem:labelwork} use only
that the completion is dominant, lies in $S_n$, and satisfies
the length identity, and that its first column has height
equal to the number of rows of the diagram of each $\tilde\tau_i$.
They therefore give $|\tilde A_i|\leq\widetilde V$, and a label
construction with at most $L\widetilde VQ_{\tilde h}(n)$ edges
costing $O(n^3+Ln^3\widetilde VQ_{\tilde h}(n))$.
Likewise, the proof of Theorem~\ref{thm:elementary} gives
\begin{equation}\label{eq:strictexpansion}
 \Sch_v(x;z)=
 \sum_{(P_1,\ldots,P_L)}
 \prod_{i=1}^L
   \varepsilon(P_i)E_{\tilde\theta_i-b_i,\tilde\theta_i}(x;Z(P_i)),
\end{equation}
summed over the filtered distinct-position paths ending at
$\tilde a_*$.

\begin{theorem}[Filtered quantum expansion]\label{thm:qexpansion}
With the same index set as \eqref{eq:strictexpansion},
\begin{equation}\label{eq:qexpansion}
 \widetilde\Sch_v(x;z)=
 \sum_{(P_1,\ldots,P_L)}
 \prod_{i=1}^L
   \varepsilon(P_i)\widetilde E_{\tilde\theta_i-b_i,\tilde\theta_i}
     (x;Z(P_i)).
\end{equation}
\end{theorem}
\begin{proof}
Apply $\mathcal Q$ to \eqref{eq:strictexpansion}.
On the left this gives $\widetilde\Sch_v(x;z)$.
Each factor in a summand is a $\mathbb Z[z]$-combination of
$e_c(X_{\tilde\theta_i})$, and the heights $\tilde\theta_i$ are
pairwise distinct.
Expanding a summand therefore gives a $\mathbb Z[z]$-combination of
standard elementary monomials.
Hence $\mathcal Q$ maps the summand to the product of the quantized
factors.
\end{proof}

Without the filters, \eqref{eq:qexpansion} is
\cite[Theorem~4.6]{samuelquantum}.
The proof shows more: for \emph{any} set of complete distinct-height
column paths, the sum of quantized weights is the quantization of
the sum of classical weights.
Thus every boundary filter justified classically is automatically
justified in quantum cohomology, without a new cancellation argument.

\subsection{The quantum Pieri step}

A \emph{quantum cover} of $w$ is a transposition $t_{ab}$, $a<b$, with
$\ell(wt_{ab})=\ell(w)+1$, or with $w(a)>w(b)$ and
$\ell(wt_{ab})=\ell(w)+1-2(b-a)$.
The latter is equivalent to every $w(c)$, $a<c<b$, lying strictly
between $w(b)$ and $w(a)$, and it carries the weight
$q_a\cdots q_{b-1}$.
The routine $\mathtt{qpieri}(w,p,k)$ of
\cite[Method~4.12]{samuelquantum} is Method~\ref{method:pieri} with
quantum covers in place of Bruhat covers, with $a\leq k<b$, with each
left position used at most once, and with the right positions $b$
weakly decreasing along a chain.
It returns triples $(w',t,q^{\delta'})$, one for every node of its
search tree.
By \cite[Lemma~4.13]{samuelquantum} it returns each endpoint once and,
for $p'\leq p\leq k$,
\begin{equation}\label{eq:qpieri}
 \widetilde\Sch_w(x;y)\widetilde E_{p',k}(x;Z)
 =\sum_{\substack{(w',t,q^{\delta'})\\ t\leq p'}}
   q^{\delta'}E_{p'-t,k-t}\bigl(Y(w,w';k);Z\bigr)
   \widetilde\Sch_{w'}(x;y).
\end{equation}
Here $t$ is the number of transpositions in the chain, and
$Y(w,w';k)$ is as in Section~\ref{sec:pieri}.
Since left positions are used once, it has $k-t$ entries.
Multiplicity-freeness gives the analogue of
Proposition~\ref{prop:pieriwork}: if all nodes lie in $S_M$ and $B$
triples are returned, the call costs $O(M^3B)$ elementary operations
and $O(MB)$ ring operations for the $q$-monomials.

\begin{method}[Quantum \alg]\label{method:qproduct}
Construct the strict completion and its label graphs.
The coefficient ring is $\mathbb Z[q][y,z]$.
Initialize $H(u,\id)=1$.
At column $i$ put $k=\tilde\theta_i$ and create a fresh table $H'$.
\begin{enumerate}
\item Group the keys of $H$ by their first component $w$.
For each such $w$ let $p_w$ be the maximum of $k-b$ over all label
edges, of length $b$, leaving labels paired with $w$.
Call $\mathtt{qpieri}(w,p_w,k)$ once.
\item For each table entry $H(w,a)$, each edge $P:a\to a'$ of length $b$,
and each returned triple $(w',t,q^{\delta'})$ with $t+b\leq k$, accumulate
\[
 H'(w',a')\mathrel{+}=
  H(w,a)\,\varepsilon(P)\,q^{\delta'}\,
  E_{k-t-b,k-t}\bigl(Y(w,w';k);Z(P)\bigr).
\]
\item Replace $H$ by $H'$.
After column $L$, return $H(w,\tilde a_*)$; its coefficient of
$q^\delta$ is $c_{u,v}^{w,\delta}(y;z)$.
\end{enumerate}
In the ordinary quantum specialization keep only $t+b=k$.
\end{method}

\begin{proposition}\label{prop:qcorrect}
Method~\ref{method:qproduct} computes \eqref{eq:qproduct}.
\end{proposition}
\begin{proof}
Multiply \eqref{eq:qexpansion} by $\widetilde\Sch_u(x;y)$ and apply
\eqref{eq:qpieri} at each column, with $p'=k-b$.
The coefficients in \eqref{eq:qpieri} involve only $q,y,z$, so they
commute with the remaining factors.
The proof of Theorem~\ref{thm:correct} applies verbatim, and no
contribution is lost since $p'\leq p_w$.
\end{proof}

\subsection{States of the quantum computation}

Write $D(w)=r_w$ for the last descent.
Put
\[
 \rho=\max(r_u,\tilde h),\qquad
 s(d)=\max\{j\geq0:j(j+1)\leq d\},\qquad
 R_q=\min\bigl(M-1,\ \rho+\min(L,s(d))\bigr),
\]
and
\[
 U_q=\min\left\{\fall{M}{R_q},\binom{d+R_q}{R_q}\right\},
 \qquad M=n+L.
\]

\begin{lemma}[Quantum descent steps]\label{lem:qdescent}
Let $t_{ab}$ be a quantum cover of $w$ with $a\leq k<b$.
If it is a Bruhat cover then $D(wt_{ab})\leq\max(D(w),k)$.
If it is length-decreasing then $b\leq D(w)+1$ and
$D(wt_{ab})\leq\max(D(w),b)$.
In particular, a single quantum cover raises $D$ by at most one.
\end{lemma}
\begin{proof}
Put $w'=wt_{ab}$ and $D=D(w)$.
The entries of $w$ after position $D$ increase.
Only positions $a$ and $b$ change, so a descent of $w'$ after
$\max(D,k)$ can only occur at $b-1$ or $b$ with $b>D$.

For a Bruhat cover, $w'(b)=w(a)<w(b)<w(b+1)$, so there is no descent at
$b$.
If $b-1=a$ there is no descent at $b-1$ because $w'(a)=w(b)>w(a)$.
Otherwise $w(b-1)<w(b)$, and the cover condition forces
$w(b-1)<w(a)=w'(b)$.

For a length-decreasing cover, suppose $b\geq D+2$.
If $a<b-1$ then $w(b-1)$ lies between $w(b)$ and $w(a)$, so
$w(b-1)>w(b)$; but $b-1>D$.
If $a=b-1$ then $a>D$ and $w(a)<w(b)$.
Either case is a contradiction, so $b\leq D+1$.
The entries after position $b$ are unchanged and increasing,
so $D(w')\leq\max(D,b)$.
\end{proof}

\begin{lemma}[Quantum state bounds]\label{lem:qstates}
Every permutation produced by a $\mathtt{qpieri}$ call in
Method~\ref{method:qproduct} lies in $S_M$, has length at most $d$,
and has no descent after $R_q$.
There are at most $U_q$ such permutations.
Every monomial $q^\delta$ that occurs involves only
$q_1,\ldots,q_{M-1}$ and has $|\delta|\leq d/2$.
\end{lemma}
\begin{proof}
\emph{Ambient size.}
Let $w\in S_N$ and $k<N$.
For $c>N$, $w(c)=c$.
Hence a first transposition with $b\geq N+2$ is not a quantum cover:
the value $N+1$ at position $N+1$ lies strictly between $w(a)$ and
$w(b)=b$, which is forbidden for a Bruhat cover, and $w(a)<w(b)$
rules out a length-decreasing cover.
Since the values $b$ weakly decrease, the call stays in $S_{N+1}$.
At column $i$ the height is $k\leq\tilde h\leq n-1$, so induction
gives $S_{n+i}\subseteq S_M$.
In particular, every quantum factor $q_a\cdots q_{b-1}$ has
$b\leq M$.

\emph{Length and $q$-degree.}
Classical covers add one to $\ell$; a length-decreasing cover $t_{ab}$
changes $\ell$ by $1-2(b-a)$ and multiplies by a monomial of degree
$b-a$.
Hence, along any chain of covers starting at $u$,
\begin{equation}\label{eq:qbalance}
 \ell(w)+2|\delta|=\ell(u)+T,
\end{equation}
where $T$ is the total number of transpositions.
Consider a sequence of pairs reaching a key $(w,a)$ at boundary $i$, with
chain lengths $t_j$ and label edge lengths $b_j$.
Then $t_j\leq\tilde\theta_j-b_j$, and the labels are length-additive,
so
\[
 T=\sum_{j\leq i}t_j\leq|\tilde\theta_{\leq i}|-\ell(a),\qquad
 |\tilde\theta_{\leq i}|=\sum_{j\leq i}\tilde\theta_j.
\]
The label $a$ lies on a complete label path to $\tilde a_*$, whose
later edges have total length at most $\sum_{j>i}\tilde\theta_j$.
Lemma~\ref{lem:strict} gives
$\ell(\tilde a_*)=|\tilde\theta|-\ell(v)$, so
$\ell(a)\geq|\tilde\theta_{\leq i}|-\ell(v)$ and $T\leq\ell(v)$.
A node produced from $w$ in column $i+1$ has $t\leq p_w=k-b$ for an
edge $a\to a'$ leaving a label paired with $w$.
Concatenating with a sequence reaching $(w,a)$ gives
$T+t\leq|\tilde\theta_{\leq i+1}|-\ell(a')\leq\ell(v)$, by the same
argument at boundary $i+1$.
By \eqref{eq:qbalance} that node has length at most $d$, and the
monomial along this sequence has degree at most $d/2$.
All monomials attached to the node come from such sequences.

\emph{Descents.}
Fix such a sequence ending at a node $w$.
All heights satisfy $k\leq\tilde h\leq\rho$, and $D(u)\leq\rho$.
By Lemma~\ref{lem:qdescent}, a cover applied to a permutation with
$D\leq\rho+j-1$ produces one with $D\leq\rho+j$, and the bound
$D\leq\rho+j$ is reached from $D=\rho+j-1$, for $j\geq1$, only by a
length-decreasing cover with $b=\rho+j$.
Since $a\leq k\leq\rho$, that cover has $q$-degree $b-a\geq j$.
If $D(w)=\rho+e$ with $e>0$, the sequence therefore contains covers
of $q$-degrees at least $1,2,\ldots,e$, so $e(e+1)/2\leq|\delta|\leq d/2$.
Furthermore, suppose a call first raises $\max(D,\rho)$ to $\rho+j$.
This happens through a cover with $b=\rho+j$.
All later covers in that call have $b\leq\rho+j$, and by
Lemma~\ref{lem:qdescent} they keep $D\leq\rho+j$.
Hence each call raises $\max(D,\rho)$ by at most one.
A sequence through column $i$ makes $i\leq L$ calls.
This proves $D(w)\leq R_q$; the bound $D(w)\leq M-1$ is automatic.

The count $U_q$ follows as in Lemma~\ref{lem:states}.
\end{proof}

\begin{proof}[Proof of Theorem~\ref{thm:quantum}]
Correctness is Proposition~\ref{prop:qcorrect}.
The cost is counted exactly as in the proof of Theorem~\ref{thm:main},
with $U,V,h,\theta$ replaced by $U_q,\widetilde V,\tilde h,\tilde\theta$,
using Lemmas~\ref{lem:strict} and \ref{lem:qstates} and the strict
label bound above.
The table has at most $U_q\widetilde V$ keys because
monomials in $q$ are part of the coefficient.
The ambient and degree assertions are in Lemma~\ref{lem:qstates}.
For the coarse bound, use $L\leq\tilde h\leq n-1$, $M\leq2n-1$,
$U_q\leq(2n)^{R_q}$, $\widetilde V\leq n^{\tilde h}$,
$Q_{\tilde h}(n)\leq2n^{\tilde h}$, and $R_q\leq\rho+L\leq2\rho$.
The final claims follow from \eqref{eq:strictheight}.
\end{proof}

\subsection{Remarks on the quantum bound}

\emph{Dependence on $L$.}
Classically $h=r_v$, so fixing code lengths suffices.
In the quantum algorithm the factor $\tilde h$ also requires $L$ to be
bounded.
This is not a weakness of the state estimate alone: the example
$\Sch_v=x_1^m$ has $r_v=1$ and $\tilde h=m$.
The difference comes from the distinct-height condition in
Theorem~\ref{thm:qexpansion}.
Both $L$ and $\tilde h$ are computed from the input code of $v$, and
$L\leq\min(\ell(v),n-1)$.
For ordinary quantum multiplication one may interchange $u$ and $v$.

\emph{The worst-case parameters are pessimistic.}
For $u=v=(m+1)12\cdots m$, $2\leq m\leq8$, the measured numbers of
distinct $u$-side states per column were $m^2+m+1$, and the label sets
had $m$ elements. The corresponding bounds $U_q$ and $\widetilde V$ grow
superpolynomially.
We do not know whether polynomial bounds in $n$, for fixed code
lengths, hold for the strict label graphs.

\emph{The implemented medium shape.}
The package's fast quantum kernels use the medium shape
$\bar\theta$ of \cite[Definition~4.1]{samuelquantum}, which allows
certain pairs of equal heights and has $\bar\theta_1\leq\tilde h$
\cite[Proposition~4.3]{samuelquantum}.
Pairs are processed by the doubled Pieri rule
\cite[Methods~4.14--4.15]{samuelquantum}.
The termwise-quantization argument also applies to these shapes.
Write $X_{k+1}=X_k\cup\{x_{k+1}\}$.
Splitting $e_c(X_{k+1})=e_c(X_k)+x_{k+1}e_{c-1}(X_k)$ and reindexing
proves
\[
 e_a(X_k)e_b(X_k)
 =\sum_{c\leq b}e_c(X_{k+1})e_{a+b-c}(X_k)
  -\sum_{c>a}e_c(X_{k+1})e_{a+b-c}(X_k).
\]
In a medium shape the column before a pair of height $k$ has height
at least $k+2$, so the heights $k,k+1$ are not used elsewhere.
Thus $\mathcal Q$ maps each filtered medium summand to the product of
its single quantized factors and its doubled factors
\cite[(4.1)]{samuelquantum}.
This gives a filtered version of \cite[Theorem~4.7]{samuelquantum}.
A doubled step must, however, retain the intermediate permutation of
its two chains, so the same accounting gives a bound with $U_q^3$ in
place of $U_q^2$.
We do not analyze the implementation's additional degree cap.

\emph{Prior work.}
Quantum Schubert polynomials \cite{fomin1997quantum}, the quantum Pieri
rule \cite{postnikov1999quantum}, and equivariant quantum Schubert calculus
\cite{mihalcea2006equivariant,lam2014quantum} give algorithms
for these coefficients.
We are not aware of a previous explicit bound for computing an entire
quantum product on complete flag varieties in terms of input
parameters.
Since the $q^0$ part of \eqref{eq:qproduct} is \eqref{eq:product},
the hardness remarks in Section~\ref{sec:prior} apply to it.

\section{Relation to the implementation and scope}\label{sec:implementation}

The package's ordinary and double kernels use the same main data flow,
as Proposition~\ref{prop:uniform} predicts:
compute the dominant code of $v^{-1}$, construct backward label graphs,
enumerate Pieri endpoints once per first-component state, and
accumulate matched contributions.
Methods~\ref{method:labels} and \ref{method:product} follow the classical
algorithm of \cite[Methods~4.9--4.11]{samuelquantum}.
In particular, the package's \texttt{kdown\_perms} restriction to positions
still holding their original value is the distinct-position condition
justified by Lemma~\ref{lem:cancellation}.
The Cauchy filter in Lemma~\ref{lem:filter} separately explains which
column endpoints may be passed to the next layer.
Neither assertion relies on empirical cancellation.

The remaining distinction concerns arithmetic evaluation.
The package's elementary-weight routine uses an unmemoized
divide-and-conquer expression construction.
The circuit implementation analyzed in Theorem~\ref{thm:main} uses
the shared recurrence
\eqref{eq:weightrec}.
For a literal recursive implementation, a conservative upper bound
is obtained by multiplying \eqref{eq:main} by
$(h+1)^{O(\log(h+1))}$.
This leaves the general $2^{O(n\log n)}$ bound and fixed-code-length
polynomiality unchanged, but is not the same operation count.

Arithmetic-circuit complexity is also not a wall-clock bound for
a computer algebra system that canonicalizes expressions, copies
subtrees, or performs automatic simplification.
Probabilistic removal of zero coefficients is not part of the
algorithm analyzed here.
Theorems~\ref{thm:main} and \ref{thm:quantum} apply to the explicitly
specified algorithms: the triple Method~\ref{method:product} and the
strict-shape quantum Method~\ref{method:qproduct}, together with their
double and ordinary specializations.
They do not cover Grothendieck kernels, the medium-shape quantum kernel
with its doubled Pieri steps (see Section~\ref{sec:quantum}),
merged-column implementations, or every released optimization.
The package's reference quantum routines use the strict shape and
single-column steps, as Method~\ref{method:qproduct} does.

As finite checks, the specified update was compared with the
ordinary Python kernel on all $616$ pairs from $S_2,S_3,S_4$.
The double update was checked on all pairs from $S_3$ at three
integer specializations and on $40$ additional pairs from $S_4$,
both against the double Python kernel and using independent
divided-difference evaluation of the Schubert polynomials on both
sides of \eqref{eq:product}.
The dominant-completion and descent assertions were also checked
through $S_6$, and the boundary counts on $528$ weak intervals
through $S_5$.
The label graphs and their coefficient alphabets were compared with
the implementation on all $148$ nonidentity inputs from $S_2$ through $S_5$.
The distinct-position enumeration was compared with
\texttt{kdown\_perms}, and with the unrestricted weighted path sum,
on $2728$ column instances through $S_5$.
The elementary recurrence was checked in all $91$ degrees
$0\leq p\leq m\leq12$, including repeated coefficient parameters.
For Corollary~\ref{cor:grass}, the rectangular completion and the
Grassmannian property of every enumerated $u$-side state were checked
on $4233$ pairs of Grassmannian permutations through $S_7$.

For the quantum section, Method~\ref{method:qproduct} was implemented
independently, with $\mathtt{qpieri}$ written from the definition of a
quantum cover and compared with \texttt{elem\_sym\_perms\_q} on every
call.
Its ordinary specialization agreed with both package quantum kernels,
the strict-shape reference and the medium-shape fast kernel, on all
$584$ pairs with $v\neq\id$ from $S_2,S_3,S_4$ and on $149$ random
pairs from $S_5$.
The triple coefficients agreed with both quantum double kernels at
random integer specializations of $y,z,q$ on all $30$ such pairs from
$S_3$ and $60$ random pairs from $S_4$.
At $z=y$ the computed coefficients satisfied
$c_{u,v}^{w,\delta}=c_{v,u}^{w,\delta}$ on $555$ pairs through $S_4$.
In all these runs the assertions of Lemma~\ref{lem:qstates} and the label
bound $|\tilde A_i|\leq\widetilde V$ were checked for every
enumerated state.
Lemma~\ref{lem:strict} was checked on all $5906$ nonidentity
permutations through $S_7$.
These checks supplement the proofs; they neither certify symbolic
identities in all ranks nor prove equivalence of all optimized routines.

\section{Further questions}

The estimates deliberately separate correctness from sharp counting.
Several improvements remain meaningful without introducing output size
as a parameter.
One is to replace $Q_h(n)$ by a smaller count exploiting the
Bruhat-cover conditions, beyond the distinct-position restriction
already justified by \cite{samuelquantum}.
For Grassmannian inputs this is the only obstruction to a $2^{O(n)}$
bound in Corollary~\ref{cor:grass}.
Another is to exploit the degree cap and the actual dominant column
heights simultaneously, instead of replacing all column heights by $h$.

In the quantum case the main question is whether the dependence on
$L$ is necessary.
The example $x_1^m$ in Section~\ref{sec:quantum} has quadratically many
$u$-side states and linearly many labels, while
Theorem~\ref{thm:quantum} allows exponentially many.
A bound polynomial in $n$ for fixed $r_u,r_v$ would require either a
quantization argument tolerating repeated column heights, such as the
doubled factors of the medium shape, or a direct count of strict label
graphs.
A sharper bound on the descent excess in Lemma~\ref{lem:qstates} would
also help; in random tests through $S_9$ it never exceeded $2$.

The fixed-code-length bound also suggests a practical division of labor.
Explicit multiplication can generate whole tables for moderate numerical
code entries, while the bounded-range lattice-point method of
Pak and Slonim is designed for compressed single-coefficient inputs.
No asymptotic crossover or universal practical superiority follows from
the two available upper bounds alone.

\bibliographystyle{acm}
\bibliography{schubmult2}
\end{document}
